\section{Analytics Methods}\label{sec:methods}
\textbf{Detection (AN-1, AN-3).} Detection runs on counts by \emph{report} week, which is what a live system sees; counts by onset week change as late reports arrive and would hide the very artefact we want to test. Three detectors run per district and disease as a grouped \texttt{applyInPandas} job. \emph{EARS C2}~\cite{ears} compares the current week with the mean and standard deviation of eight earlier weeks separated by a two-week guard band, with the standard deviation floored at the Poisson value. The \emph{seasonal detector} (Farrington-style~\cite{farr}) builds the baseline from the same weeks $\pm3$ of up to three previous years, trims the highest 20\% of reference values so past outbreaks do not inflate it, rescales by the recent year-on-year level ratio (clipped to 0.5 to 2) and uses a quasi-Poisson dispersion. \emph{CUSUM} accumulates the standardised seasonal residuals with a reference value of 1 and a decay of 0.8 so old evidence fades. An alert also requires at least five cases.

\textbf{Reporting-artefact guard.} A facility that goes silent and then back-fills produces one week in which a large share of the district's reports are more than 14 days old. Weeks with a late share of at least 15\% and at least five late reports are flagged as artefacts, excluded from every baseline and unable to raise alerts. Facility-level silence (a facility with a trailing mean of at least three reports per week that sends none) is also computed and stored as report completeness.

\textbf{Forecasting (AN-4).} For each district, disease and horizon $h\in\{1,\dots,4\}$ a negative-binomial regression with log link predicts onset-week cases. Predictors use only data that are mature at the forecast origin, namely weeks at least four weeks old, because reports arrive up to three weeks late and a feature built from incomplete recent weeks would differ between training and use. They are the log mean of the last two mature weeks, the log mean of the last eight mature weeks, two annual harmonics, and weather: mean rainfall and temperature in the window four to eight weeks before the target week (known at the origin for $h\le4$) and the log of recent rainfall. The dispersion is estimated from a Poisson fit by the method of moments, and prediction quantiles come from the negative-binomial distribution. Series with a mean below 0.8 use an intercept and annual cycle only. All series are trained in parallel by Spark. Dengue, diarrhoea and malaria are backtested; ILI is forecast operationally for the bed-demand estimate.

\textbf{Demand conversion (AN-6).} Forecast cases over the next four weeks become resource needs with the assumptions in Table~\ref{tab:assume}: IV saline per admission (dengue 5 bags, diarrhoea 4), NS1 kits per dengue case (1.2), ORS sachets per diarrhoea case (5), antimalarial courses per malaria case (1) and length of stay for beds. District admission rates are estimated from mature weeks of the last year and shrunk towards the national rate with a prior strength of 50 cases. The upper demand sums the weekly 90\% quantiles, a deliberately conservative bound that treats weeks as perfectly correlated (the sum of quantiles is above the quantile of the sum), and the lower and mean demands are computed alike.
\begin{table}[t]\caption{Conversion assumptions (illustrative; one versioned file).}\label{tab:assume}\centering\small
\begin{tabular}{@{}lrrl@{}}\toprule Resource & Cost (USD) & Volume & Driver \\\midrule
IV saline bag & 1.20 & 1.00 & admissions \\ ORS sachet & 0.08 & 0.02 & diarrhoea cases \\
NS1 kit & 2.50 & 0.05 & dengue cases \\ Antimalarial course & 1.50 & 0.01 & malaria cases \\\bottomrule\end{tabular}\end{table}

\textbf{Allocation (AN-7, AN-8).} With districts $d$ and resources $r$, need $n_{dr}=\max(0,\hat y_{dr}-s_{dr})$ (upper demand minus stock), allocation $x_{dr}\ge0$ and unmet demand $u_{dr}\ge0$, the programme minimises $\sum w_dc_ru_{dr}$ (plus a tiny spending term that avoids pointless allocation) subject to $u_{dr}+x_{dr}\ge n_{dr}$, a national budget $\sum c_rx_{dr}\le B$, an equity floor $x_{dr}\ge\alpha p_d\sum_{d'}x_{d'r}$ and a transport limit $\sum_{d\in k}\sum_rv_rx_{dr}\le T_k$ for each division $k$. Valuing unmet demand in money makes resources comparable. The severity weight $w_d=1+(1-b_d/b_{\max})$ rises as hospital beds per 100\,000 people ($b_d$) fall. The transport limit is 1.5 times the division's population share of the national volume of need. The equity floor $\alpha$ is a policy dial: at $\alpha=0$ the model sends supplies wherever they avert the most weighted harm and may leave low-risk districts with nothing; at $\alpha=0.4$ every district is guaranteed at least 40\% of its population share of every resource, whatever its forecast. The baseline is a population-proportional split of the same budget. The programme is solved with SciPy HiGHS.

\textbf{Seasonal calendar (AN-9).} For every division and disease the weekly incidence per 100\,000 is decomposed by STL (period 52, robust). The seasonal component averaged by epi-week defines the high season as the contiguous window around the peak where it exceeds its mean by half a standard deviation; supplies should be pre-positioned six weeks and a prevention campaign started eight weeks before the season begins (stated lead-time assumptions). A season is called marked only if the seasonal strength is at least 0.3.
